%=============================================================================!
% 16.4  REACTIONS AND THEIR RATES                                             !
%=============================================================================!
\section{Reactions and their rates}
\label{sec:ob-reactions}

Once molecules have been identified, a reaction is a change in which
molecules a group of atoms belongs to. We record the reactants and
products, then turn repeated events into rates. The uncertainty depends
on the number of independent events, rather than the number of saved
frames. We also measure the cost of doing this analysis during a run.

\subsection*{Reactions from changes of species}

Between two frames, each atom belongs to a molecule before and to a
molecule after. Joining each old molecule to each new one that shares an
atom with it gives a graph whose nodes are molecules, and its connected
components are the reactions. A molecule that kept its atoms and its
bonds is one old node joined to one new node with the same hash, and is
passed over; every other component is a reaction, its old molecules the
reactants and its new ones the products. A ring that opens is a
component of one old and one new molecule with different hashes, C$_6$
$\to$ C$_6'$; two molecules that exchange an atom, AB + B $\to$ B + AB, are
a component of two old and two new molecules with the same names.
\texttt{mdlab.bonds.ReactionLog} does this frame by frame, recording the
time, the reactants and the products of each reaction, and the counts of
each species.

In the A-B gas at \SI{2000}{\kelvin} the 250 molecules AB lose partners
at first (\cref{fig:ob-bonds}b), and after the first \SI{10}{\pico
\second} the gas holds on average 139.3 AB, 104.1 free A and 103.9 free B,
with 2.2 AB$_2$, 1.9 A$_2$B and a few larger complexes.
Over \SI{100}{\pico\second} the log holds 57 kinds of reaction. The most
frequent, counting every event with the band and no least lifetime, are
\begin{center}
\begin{tabular}{@{}lr@{\qquad}lr@{}}
AB $\to$ A + B & 4517 & A + B $\to$ AB & 4421\\
AB + B $\to$ AB$_2$ & 1618 & AB$_2$ $\to$ AB + B & 1602\\
A + AB $\to$ A$_2$B & 1556 & A$_2$B $\to$ A + AB & 1566\\
AB + AB $\to$ A$_2$B$_2$ & 301 & A$_2$B$_2$ $\to$ AB + AB & 298
\end{tabular}
\end{center}
a reaction network: a graph whose nodes are the species and the
reactions, each reaction joined to its reactants and its products and
carrying its count. Each reaction and its reverse occur
nearly equally often, as every process in a gas at equilibrium is
balanced by its reverse: AB $\to$ A + B outnumbers its reverse by 96, about
the 111 net breaks by which the molecules fell from 250 to 139.3 at the
start, and the others differ by less than the $\sqrt n$ of the toolbox
below. The complexes are brief. In a steady state the complexes present
at any moment are those formed within one lifetime before it, so their
number is the rate at which they form times the lifetime. AB$_2$ forms
16.8 times per picosecond and 2.18 are present on average, so each lasts
$2.18/16.8 = \SI{0.130}{\pico\second}$; A$_2$B lasts
\SI{0.120}{\pico\second}. They are passing
encounters of the kind the least lifetime of \cref{sec:ob-bonds} removes.

\begin{toolbox}[Rates from counts, and the Poisson distribution]
Events that happen independently, at a steady rate $\nu$ for each of
many molecules, are counted over an exposure $X$, the sum over the
molecules of the time each was able to react. Cut into $J$ small pieces,
the exposure has in each an event with the small probability $p = \nu
X/J$ and none otherwise, independently, so the count is a sum of $J$
independent numbers each 0 or 1. Each such number has the mean $p$, and
so does its square, which is the same 0 or 1, so its variance is $p -
p^2 = p(1 - p)$. The count then has the mean $Jp = \nu X$ and, since the
variances of independent numbers add (\cref{sec:cv-mean}), the variance
$Jp(1 - p)$, which tends to $\nu X$ as the pieces shrink. The rate is
estimated as $n/X$, with the standard error $\sqrt n/X$, a fraction
$1/\sqrt n$ of itself: a rate from 100 events is known to a tenth, from
10\,000 to a hundredth. The chance of $n$ events is the number of ways of
choosing $n$ of the $J$ pieces, $J(J - 1)\cdots(J - n + 1)/n!$, times
$p^n(1 - p)^{J - n}$, which with $p = \nu X/J$ is
\begin{equation*}
   \frac{J(J - 1)\cdots(J - n + 1)}{J^n}\,\frac{(\nu X)^n}{n!}\,
   \Bigl(1 - \frac{\nu X}{J}\Bigr)^{J - n} .
\end{equation*}
As $J$ grows with $n$ fixed, the first fraction, a product of $n$
factors each tending to 1, tends to 1, and the logarithm of the last
factor, $(J - n)\ln(1 - \nu X/J)$, tends to $-\nu X$, since $\ln(1 + y)$
is $y$ plus terms of order $y^2$ (\cref{sec:th-langevin}). The chance
tends to the Poisson distribution
\begin{equation*}
   \mathcal{P}(n) = \frac{(\nu X)^n}{n!}\,\mathrm{e}^{-\nu X} ,
\end{equation*}
whose mean and variance are both $\nu X$, the limits of the count's.
\end{toolbox}

\noindent Counted events of a run can be correlated, a bond that breaks
and re-forms twice counting thrice; the least lifetime of
\cref{sec:ob-bonds} removes most of these. At \SI{2000}{\kelvin}, the bonds
broken in 89 windows of \SI{1}{\pico\second} number 33.82 on average with
variance 36.24, a ratio of 1.07, within the 0.15 by which that ratio
scatters over 89 windows of a Poisson count, and their histogram lies
close to the Poisson distribution of the same mean
(\cref{fig:ob-rates}b).

\subsection*{Rates against temperature}

With the band and a least lifetime of \SI{0.5}{\pico\second}, after the
first \SI{10}{\pico\second}, the rate of breaking is the number of bonds
broken divided by the bond-picoseconds of exposure, and the rate constant
of forming the number formed divided by the exposure of free pairs,
$\int n_{\mathrm A}n_{\mathrm B}/V\,\dd t$. At \SI{1500}{\kelvin} 1828 bonds
break over \num{18168} bond-picoseconds, $k_{\mathrm b} = (0.1006 \pm
0.0024)$\,\si{\per\pico\second}, and at \SI{2500}{\kelvin} 3182 over 9500,
$(0.3350 \pm 0.0059)$\,\si{\per\pico\second}; the rate constant of forming
falls over the same range, from $(1911 \pm 45)$ to $(415 \pm
7)$\,\si{\angstrom\cubed\per\pico\second} (\cref{fig:ob-rates}a). A rate
that rises as $\mathrm{e}^{-E_{\mathrm a}/\kB T}$ has a logarithm that
falls on a straight line of slope $-E_{\mathrm a}$ against $1/\kB T$, the
Arrhenius form. The error of the logarithm of a rate from $n$ events is
its fractional error, $1/\sqrt n$, so each point is weighted in the fit
by $\sqrt n$, the misses being measured in units of their errors as in
\cref{sec:cv-size}: breaking has $E_{\mathrm a} = \SI{0.389}{\electronvolt}$
and forming $-\SI{0.478}{\electronvolt}$, a rate that falls as the
temperature rises.

Neither slope is the depth of the well, \SI{1}{\electronvolt}, and they
need not be: each rate depends on how the buffer gas supplies and removes
energy, as Chapter~13 found for hops. Their ratio does not. At
equilibrium forming and breaking balance, $k_{\mathrm f}n_{\mathrm A}
n_{\mathrm B}/V = k_{\mathrm b}n_{\mathrm{AB}}$, so $k_{\mathrm f}/k_{\mathrm
b}$ is the equilibrium constant $K_{\mathrm{eq}} = n_{\mathrm{AB}}V/n_{\mathrm
A}n_{\mathrm B}$ of the counts: 18\,994 against 20\,583
\si{\angstrom\cubed} at \SI{1500}{\kelvin}, 3406 against 3266 at 2000 and
1237 against 1177 at \SI{2500}{\kelvin}, within 8\%. The two are not
the same average. With nearly as many bonds formed as broken,
$k_{\mathrm f}/k_{\mathrm b}$ is the bond-picoseconds over the exposure of
free pairs, 19\,162, 3416 and \SI{1235}{\angstrom\cubed}, a ratio of time
averages that counts the bonds within complexes too, while $K_{\mathrm{eq}}$
averages the ratio of the counts for the molecules AB alone. Since
$\ln K_{\mathrm{eq}} = \ln k_{\mathrm f} - \ln k_{\mathrm b}$, its slope is the
difference of theirs, $0.478 - (-0.389) = \SI{0.867}{\electronvolt}$, and
the counts give \SI{0.912}{\electronvolt}.

That constant can be predicted for a dilute gas of A and B whose only
interaction is the pair potential, with a pair bonded when closer than
$r_{\mathrm{on}}$. Arrangements with $n_{\mathrm{AB}}$ bonded pairs and
with one more, made by joining a free A to a free B, are compared by the
Boltzmann weights of \cref{sec:sm-canonical} integrated over the
positions; the momenta, the same in both, cancel. A free pair, each atom
anywhere, has the weight $V^2$; a bonded pair, its centre anywhere and
its separation within $r_{\mathrm{on}}$, has $V\int_0^{r_{\mathrm{on}}}4\pi
r^2\mathrm{e}^{-\beta\varphi(r)}\,\dd r$. The A and the B to join can be
chosen in $n_{\mathrm A}n_{\mathrm B}$ ways, and each result is reached
from $n_{\mathrm{AB}} + 1$ starting arrangements, one for each of its
bonds, so one more bond is more probable than not by the factor
\begin{equation*}
   \frac{n_{\mathrm A}n_{\mathrm B}}{n_{\mathrm{AB}} + 1}\,\frac1V
   \int_0^{r_{\mathrm{on}}}4\pi r^2\mathrm{e}^{-\beta\varphi(r)}\,\dd r .
\end{equation*}
At the most probable $n_{\mathrm{AB}}$, large enough that $n_{\mathrm{AB}} +
1$ differs little from $n_{\mathrm{AB}}$, the factor is 1, so
$K_{\mathrm{eq}}$ equals the integral, the mass action law. The integral falls
sixteen-fold from \num{15756} to \SI{987}{\angstrom\cubed} between 1500
and \SI{2500}{\kelvin}, and the constant from the single cutoff at
$r_{\mathrm{on}}$, with the standard error of nine blocks, follows it at
1.00, 0.89, 0.92, 0.95 and 0.91 of its value (\cref{fig:ob-bonds}c); the
band, which also keeps pairs between $r_{\mathrm{on}}$ and $r_{\mathrm{off}}$
that were once closer, gives 1.13 to 1.31. The integral treats the free
atoms as an ideal gas, leaving out the attraction between unlike atoms
and the repulsion between like ones beyond $r_{\mathrm{on}}$.

\subsection*{Carbon from the melt}

Carbon's bonds come and go in the melt. 216 atoms at
\SI{2.0}{\gram\per\centi\metre\cubed}, with the potential of
Tersoff\cite{Tersoff1989} through ASE's calculator, the parameters those
that ship with ASE, are held at \SI{6000}{\kelvin} for
\SI{2}{\pico\second} by CSVR with $\tau_{\mathrm T} = \SI{50}{\femto
\second}$, and the thermostat's target is then lowered steadily to
\SI{300}{\kelvin}, over \SI{10}{\pico\second} in one run and
\SI{40}{\pico\second} in another, with steps of \SI{0.5}{\femto\second}.
The model is carried by \texttt{io.AseModel}, which calls any ASE
calculator with the positions and returns the energy and the forces as an
\texttt{mdlab} model does. The default band, 1.2 to 1.4 times twice the
covalent radius of carbon, \num{1.824} to \SI{2.128}{\angstrom}, brackets
the first minimum of $g(r)$ in the melt, at \SI{1.963}{\angstrom}, and
after the quench lies, all but its first \SI{0.03}{\angstrom}, in the gap
from 1.85 to \SI{2.15}{\angstrom} where $g(r)$ is zero
(\cref{fig:ob-carbon}a), so the criterion counts the first shell of
neighbours and nothing beyond. In the melt, between 0.5 and
\SI{2}{\pico\second}, 0.764 of the atoms have three bonds, 0.161 two and
0.071 four, and the bonds form and break continually: after the first
frame there are 1870 events over the faster run and 3479 over the slower,
1221 of each in the first \SI{2}{\pico\second} at \SI{6000}{\kelvin}. As the melt cools the three-fold atoms grow at the expense of
the others, and the last bond forms or breaks when the thermostat's
target has fallen to \SI{2335}{\kelvin} in the faster run and
\SI{2432}{\kelvin} in the slower (\cref{fig:ob-carbon}b). At \SI{300}{\kelvin} the faster quench leaves
0.852 of the atoms with three bonds, 0.065 with two and 0.083 with four;
the slower leaves 0.907, 0.042 and 0.051. Every atom is in one molecule,
the solid itself. These are single runs, so the difference between the
two quenches has no error bar yet; the trend, fewer defects for slower
cooling, is what a repeat would test.

\subsection*{Keeping up with the run}

An analysis can read frames from a file after the run, or watch the run
as it goes. \texttt{io.read\_frames} reads any trajectory file ASE can
read, extended XYZ or ASE's own, converting velocities to
\si{\angstrom\per\femto\second}. Watching needs no file: \texttt{md.run}
and \texttt{thermostats.run} accept observers, functions called at every
recorded frame with the time, the positions and the velocities, and the
pairs that the force evaluation has just found are at hand in
\texttt{PairModel.last\_pairs}, so an observer can feed them to a
\texttt{ReactionLog} without searching for neighbours again. The gas runs
above did so every \SI{10}{\femto\second}.

Such analysis must cost little next to the forces. \Cref{fig:ob-timing}
times one frame's analysis, bonds, molecules, hashes, formulas and
reactions, against one force evaluation of a cheap model, liquid argon's Lennard-Jones forces with the neighbour list already
built, from 864 to \num{97556} atoms. The bonds are taken between atoms
closer than \SI{4.0}{\angstrom}, kept to \SI{4.4}{\angstrom}, 4.4 per atom, a
denser graph than most molecules give. With the two loops of
\cref{sec:ob-bonds} and the rounds of the hash compiled by Numba, and the
pairs taken from the force evaluation, the analysis costs about 0.15 of a
force evaluation at 864 atoms and about 0.10 from \num{10976} atoms up,
in each of two sets of timings. From 864 to \num{97556} atoms, 113 times as
many, both times grow roughly in proportion to $N$, the force
evaluation's somewhat faster. Analysed at a
frame saved every tenth step, the bonds cost about 1\% of the run, and
with any learned potential, whose forces cost far more than these, much
less.

\begin{figure}[htb]
\centering
\includegraphics{ch16_observables/rates}
\caption{The A-B gas. (a) The rate of breaking per bond and picosecond
(circles) and the rate constant of forming in \SI{e3}{\angstrom\cubed\per
\pico\second} (squares), counting events of the band that lived at least
\SI{0.5}{\pico\second}, with the error $\sqrt n$ of each count, against
$1/\kB T$, with straight lines fitted to their logarithms. (b) At
\SI{2000}{\kelvin}, the fraction of 89 windows of \SI{1}{\pico\second} in
which a given number of bonds broke (bars), against the Poisson
distribution of the same mean (points).}
\label{fig:ob-rates}
\end{figure}

\begin{figure}[htb]
\centering
\includegraphics{ch16_observables/carbon}
\caption{216 carbon atoms at \SI{2.0}{\gram\per\centi\metre\cubed} with
Tersoff's potential, quenched from \SI{6000}{\kelvin}. (a) $g(r)$ of the
melt and of the solid after the slow quench, with $r_{\mathrm{on}}$ and
$r_{\mathrm{off}}$ of the default band (dotted). (b) The fractions of atoms
with two, three and four bonds against the thermostat's target, for the
quench over \SI{10}{\pico\second} (thin) and over \SI{40}{\pico\second}
(thick).}
\label{fig:ob-carbon}
\end{figure}

\begin{figure}[htb]
\centering
\includegraphics{ch16_observables/timing}
\caption{(a) The time of one force evaluation of liquid argon's model
and of one frame's bond analysis, against the number of atoms, each
relative to the force evaluation of 864 atoms, the least of three tries.
(b) Their ratio, with the tenth that the analysis is meant to stay within
(dotted).}
\label{fig:ob-timing}
\end{figure}

\begin{keyidea}
A reaction is a change of the molecules that some atoms belong to, found
as a connected group of old and new molecules. Rates are counts over
exposures, with the error $\sqrt n$ of a Poisson count; the ratio of the
rate constants of a reaction and its reverse is the equilibrium constant,
which statistical mechanics predicts. Observers analyse the run as it goes,
reusing the force evaluation's pairs, at about a tenth of its cost.
\end{keyidea}

\notebook{16}{count reactions in the gas and put error bars on their rates}

\begin{selfcheck}
A rate is estimated from 25 counted events. To what fraction is it known,
and how many events would halve that?
\end{selfcheck}
